\section{Уравнение теплопереноса с термомеханической генерацией}

Уравнение теплопереноса с учетом термомеханической генерации записывается как:

\begin{equation}
k \nabla^2 T - C \frac{\partial T}{\partial t} = \gamma T_0 \frac{\partial \epsilon_{kk}}{\partial t},
\label{eq:heat}
\end{equation}

где \(k\) — коэффициент теплопроводности, \(C\) — объёмная теплоёмкость, \(\gamma\) — коэффициент термоупругого связывания, \(T_0\) — начальная температура, а \(\epsilon_{kk}\) — след тензора деформаций.

\subsection{2. Левая часть — классический теплоперенос}

\subsubsection{Член \( k \nabla^2 T \)}

\begin{itemize}
    \item \(k\) — коэффициент теплопроводности
    \item \(\nabla^2 T\) — лапласиан температуры (пространственная кривизна поля температуры)
\end{itemize}

\textbf{Физический смысл:}  

Описывает распространение тепла в пространстве за счёт теплопроводности (закон Фурье).

\textbf{Коэффициент теплопроводности} — это физическая величина, которая показывает, \textbf{насколько хорошо материал проводит тепло}, измеряется в \(\text{Вт/(м·К)}\).

\(\nabla^2 T\) читается как \textit{«лапласиан температуры»}:

\[
\nabla^2 T = \frac{\partial^2 T}{\partial x^2} + \frac{\partial^2 T}{\partial y^2} + \frac{\partial^2 T}{\partial z^2}.
\]

Он показывает, \textbf{как температура в данной точке отличается от температур вокруг неё}, другими словами — \textbf{насколько температура в точке «выгнута» по сравнению с соседними точками}.

\subsubsection{Член \(- C \frac{\partial T}{\partial t}\)}

\begin{itemize}
    \item \(C\) — объёмная теплоёмкость материала \(C = \rho c_p\)
    \item \(\frac{\partial T}{\partial t}\) — скорость изменения температуры во времени
\end{itemize}

\textbf{Физический смысл:}  

Описывает накопление (или расход) тепловой энергии в материале.  

Обычная теплоёмкость показывает, сколько энергии нужно добавить, чтобы повысить температуру вещества на 1 градус.  

Объёмная теплоёмкость показывает, сколько тепла нужно, чтобы нагреть 1 кубический метр материала на 1 градус.

\subsection{3. Правая часть — термомеханическая генерация тепла}

\subsubsection{Член \( \gamma T_0 \frac{\partial \epsilon_{kk}}{\partial t} \)}

\begin{itemize}
    \item \(\gamma\) — коэффициент термоупругого связывания (связан с коэффициентом теплового расширения и модулем объёмного сжатия)
    \item \(T_0\) — опорная (начальная) температура
    \item \(\epsilon_{kk} = \epsilon_{11} + \epsilon_{22} + \epsilon_{33}\) — \textbf{след тензора деформаций}, показывает объёмную деформацию
    \item \(\frac{\partial \epsilon_{kk}}{\partial t}\) — скорость изменения объёма
\end{itemize}

\textbf{Физический смысл:}  

\(\gamma\) показывает, \textbf{как изменение температуры связано с изменением объёма или давления материала} в термоупругой модели.  

То есть, когда материал нагревается:

\begin{itemize}
    \item Он \textbf{расширяется} (объём увеличивается)
    \item Давление может изменяться, если материал ограничен
\end{itemize}

\(\gamma\) связывает тепловое и механическое поведение.  

Связь с теплоёмкостями:

\[
\gamma = C_p /C_v,
\]

где:  
\begin{itemize}
    \item \(C_p\) — теплоёмкость при \textbf{постоянном давлении}
    \item \(C_v\) — теплоёмкость при \textbf{постоянном объёме}
\end{itemize}

\subsection{1️⃣ Что такое \(\epsilon\) (эпсилон)}

\(\epsilon\) — это \textbf{тензор деформаций}, матрица, которая показывает, \textbf{как материал растягивается или сжимается в разных направлениях}.

В трёхмерном пространстве тензор деформаций записывают как \(3 \times 3\) матрицу:

\[
\epsilon =
\begin{pmatrix}
\epsilon_{11} & \epsilon_{12} & \epsilon_{13} \\
\epsilon_{21} & \epsilon_{22} & \epsilon_{23} \\
\epsilon_{31} & \epsilon_{32} & \epsilon_{33}
\end{pmatrix}.
\]

\begin{itemize}
    \item Диагональные элементы (\(\epsilon_{11}, \epsilon_{22}, \epsilon_{33}\)) показывают \textbf{растяжение или сжатие вдоль осей x, y, z}
    \item Вне диагональные элементы (\(\epsilon_{12}, \epsilon_{13}, \dots\)) показывают \textbf{сдвиги}
\end{itemize}

\subsection{2️⃣ Что такое след тензора \(\epsilon_{kk}\)}

\[
\epsilon_{kk} = \epsilon_{11} + \epsilon_{22} + \epsilon_{33}.
\]

\begin{itemize}
    \item \textbf{Сумма диагональных элементов} тензора деформаций
    \item Называется \textbf{след тензора}
    \item Показывает \textbf{объёмную деформацию материала}, то есть \textbf{насколько изменился объём тела}
\end{itemize}

\begin{itemize}
    \item Если \(\epsilon_{kk} > 0\) → тело \textbf{расширяется}
    \item Если \(\epsilon_{kk} < 0\) → тело \textbf{сжимается}
    \item Если \(\epsilon_{kk} = 0\) → \textbf{объём не меняется}, только форма
\end{itemize}

\subsubsection{3️⃣ Связь с объёмной деформацией}

Объёмная деформация \(\Delta V / V\) примерно равна \textbf{следу тензора деформаций}:

\[
\frac{\Delta V}{V} \approx \epsilon_{kk} = \epsilon_{11} + \epsilon_{22} + \epsilon_{33}.
\]

То есть \textbf{след тензора показывает, как изменился объём тела из-за растяжений и сжатий по всем направлениям}.

\textbf{Физический смысл:}  

При \textbf{сжатии или расширении} тела часть механической энергии превращается в тепловую (или наоборот).  

Это эффект \textbf{термоупругого нагрева/охлаждения}.

\begin{itemize}
    \item при сжатии (\(\dot{\epsilon}_{kk} < 0\)) → нагрев
    \item при расширении (\(\dot{\epsilon}_{kk} > 0\)) → охлаждение
\end{itemize}

\subsection{4. Физическая интерпретация целиком}

\begin{quote}
Изменение температуры в теле определяется:
\begin{itemize}
    \item теплопроводностью,
    \item теплоёмкостью,
    \item и дополнительным источником тепла, возникающим из-за изменения объёма при деформации.
\end{itemize}
\end{quote}

Это уравнение характерно для:

\begin{itemize}
    \item термоупругости,
    \item динамики твёрдых тел,
    \item акусто-термических эффектов,
    \item анализа нагрева при вибрациях и ударных нагрузках.
\end{itemize}

\subsection{5. Связь с классическим уравнением теплопроводности}

Если механических деформаций нет:

\[
\frac{\partial \epsilon_{kk}}{\partial t} = 0,
\]

то уравнение \eqref{eq:heat} сводится к стандартному:

\[
C \frac{\partial T}{\partial t} = k \nabla^2 T.
\]
